% Avsnitt 7.6, ur lectures/L07/README.md.

\section{Sammanfattning}
\label{c7:sec:review}

\needspace{6\baselineskip}
\subsection*{Det här ska du kunna nu}
\begin{itemize}
  \item Förklara vad som skiljer en timer från en vanlig räknare, och varför skillnaden är en
    jämförelse.
  \item Bygga en timer som grindnät: en likhetsjämförare av en XNOR per bit, och den nollställning
    som får den att upprepa i stället för att slå över.
  \item Implementera en timer i VHDL som en återanvändbar modul, omvandla en period i sekunder till
    ett \code{TICK_COUNT}, och säga varför \code{timeout} är en puls på en klockcykel snarare än en
    nivå.
  \item Komponera en design av färdiga moduler i stället för att skriva om dem, och säga varför en
    knapp som styr en timer ändå måste passera en synkroniserare först.
\end{itemize}

\needspace{6\baselineskip}
\subsection*{Frågor att testa dig själv med}
\begin{enumerate}
  \item I vilken mening är en timer ”bara en räknare med en jämförelse”?
  \item Varför måste den interna räknaren nollställas till \code{0} när den når sitt mål, i stället
    för att lämnas att räkna vidare?
  \item Varför är \code{timeout} en puls på en enda cykel snarare än en nivå som ligger kvar hög?
  \item I CircuitVerse-versionen innebär ett ändrat målvärde att jämföraren måste kopplas om. Vad
    ersätter den omkopplingen i VHDL-versionen, och när låses dess värde fast?
  \item Varför är det fortfarande nödvändigt att köra en knapp genom en synkroniserare innan dess
    flank används för att aktivera en timer, trots att knappen inte har med räkningen att göra?
  \item \code{walking_led} skriver nästan ingen egen logik. Nämn de tre moduler den komponerar och
    vilket kapitel var och en kom från.
\end{enumerate}

\needspace{6\baselineskip}
\subsection*{Blicken framåt}
\Chapref{c8:ch} går vidare med:
\begin{itemize}
  \item Tillståndsmaskiner, bokens sista byggblock, först konstruerade för hand: tillståndsdiagram,
    tillståndstabell, Karnaughhärledd logik och ett grindnät i CircuitVerse, precis som det här
    kapitlets timer ritades innan den skrevs.
  \item Därefter i VHDL, där en uppräkningstyp och ett \code{case} ersätter tillståndstabellen och
    Karnaughdiagrammen helt.
  \item En timer är i sig en tillståndsmaskin med två tillstånd och ett enda syfte.
    \Chapref{c8:ch} generaliserar det till kretsar med många tillstånd och många övergångar mellan
    dem.
\end{itemize}
